% =====================================================================
% Section 2: Data Resource: CMCR-Colloquial
% Overleaf 可直接粘贴。要求：LNCS (llncs.cls) + splncs04 参考文献样式。
% 粘贴位置：\section{Introduction}...之后、System Architecture 之前。
% 注意：
%   1) \ref{sec:architecture} 需与你的 System Architecture 章节标签一致，
%      若该章标签不同请全局替换；
%   2) \cite{cmcr} 的条目须在 reference.bib 中（@misc，见文末注释），勿在
%      main.tex 手写 thebibliography——否则与 \bibliography 生成的列表重复；
%   3) release 脚注中的数据集 URL 为占位，开源仓库就绪后替换。
% =====================================================================

\section{Data Resource: CMCR-Colloquial}
\label{sec:data}

We construct and release CMCR-Colloquial\footnote{Dataset URL will be provided upon release.}, a Chinese colloquial question--answering dataset for medical triage, built upon peer-reviewed clinical case reports from the Chinese Clinical Case Repository (CMCR)~\cite{cmcr}. Cases are collected under a two-tier syndrome scheme comprising two core triage complaints (fever and abdominal pain) and four red-flag (emergent) syndromes (chest pain, dyspnea, altered consciousness, and severe headache). The raw candidate pool contains 1{,}543 cases. Each report is parsed from JATS XML into a structured record with seven standardized clinical slots covering patient information, auxiliary examinations, diagnosis and differential diagnosis, treatment course, clinical outcome and follow-up, and clinical discussion, while the original section texts are retained for provenance; automated slot--source consistency checks yield 1{,}341 structured cases.

Each structured case is further decomposed offline into a set of atomic clinical propositions, which are query-agnostic and precomputed: 43{,}271 propositions in total (32.3 per case on average), of which 9{,}968 (23.0\%) are red-flag propositions. Each proposition carries a category label (core symptoms and triggers, medical history, positive findings, red-flag signs, or treatment response), bilingual Chinese--English formulations, and four groups of discriminative keywords (mention, affirm, negate, and uncertain). This proposition repository serves as the precomputed foundation for the dual-question fact verification described in Sect.~\ref{sec:architecture}.

On top of the structured cases and propositions, we generate colloquial QA pairs through a four-stage reverse-construction pipeline: (i)~loading the structured cases and propositions; (ii)~generating first-person, colloquial patient consultations addressed to an AI health assistant, using only the chief-complaint and history-of-present-illness slots; (iii)~generating lay-reader-oriented triage outputs, including the recommended department, urgency level, and suggested examinations, from the differential diagnosis, treatment course, and discussion sections; and (iv)~a state-machine-managed verification workflow. Both consultations and triage outputs are generated by the LLM DeepSeek-V4.1-Flash. Red-flag cases must embed everyday phrasings of red-flag symptoms in the patient narrative (e.g., chest pain described as ``my chest feels like a heavy stone pressing down and I can hardly breathe'') and must contain no self-diagnosis; such samples are specifically designed to exercise the emergency-escalation behavior of the system. Five-dimensional automatic quality control is applied on both sides: the input side assesses colloquial tone, masking of professional terminology, sufficiency of core clinical facts, fidelity of diagnostic direction, and narrative naturalness, whereas the output side assesses conclusion fidelity, completeness of supporting evidence, accessibility of terminology, expressive naturalness, and actionability of triage advice. Each dimension is scored up to 2 points (10 in total); scores below 8 trigger regeneration, and a score of 0 on clinical-fact sufficiency or diagnostic-direction fidelity leads to immediate rejection. The resulting dataset contains 1{,}339 complete QA pairs (fever 515, abdominal pain 450, chest pain 152, dyspnea 156, severe headache 36, and altered consciousness 30), of which 374 (27.9\%) are red-flag samples.

% =====================================================================
% 参考文献条目（并入 coauthor 的 reference.bib；勿用手写 thebibliography，
% 否则会与 \bibliography{reference} 生成的列表重复出现两个 References）
% =====================================================================
% @misc{cmcr,
%   title        = {Chinese Clinical Case Repository ({CMCR})},
%   howpublished = {\url{https://cmcr.yiigle.com/}},
%   note         = {Last accessed 2026/10/26}
% }
